\section{Introduction}
\label{sec:intro}

Language-model assistants are moving into spreadsheets: a user types a request, the assistant writes a formula, and the formula lands in a cell that other cells may depend on. When it is wrong it fails in one of two ways. It fails \emph{loudly} when a hallucinated function or sheet name produces \code{\#NAME?} or \code{\#REF!}, and \emph{silently} when a reference to an empty column, a criterion that never occurs in its column, or a sum over text evaluates to a plausible number that then enters the workbook.

Research on spreadsheet language models mostly asks whether the \emph{final answer} is right~\cite{zhao2024nl2formula,li2023sheetcopilot,ma2024spreadsheetbench,zhu2026spreadsheetbench2}. Between the model and the workbook sits a layer that has received less attention: the checks, if any, that run on a candidate before it is committed. Those checks have an unusual advantage, because the workbook is open. A candidate can be tested against facts that need no model call: which sheets exist, which columns hold data and of what type, which values occur in a column, which cells already depend on which.

We ask four questions. \textbf{RQ1:} how much of the failure of an LLM spreadsheet agent can be stopped by deterministic checks that read only the user's workbook, and where do such checks stop? \textbf{RQ2:} does the answer carry over from a benchmark to formulas people actually wrote? \textbf{RQ3:} does verification-guided repair help real models, and which? \textbf{RQ4:} does repair make failures safer, or only less visible? The claim is narrower than ``spreadsheet verification'': it is deterministic verification of an LLM-generated formula against workbook state \emph{before the write}, evaluated with fault labels from an independent engine, an explicit distinction between loud and silent failures, and repair guided by the verifier. We study an open-source Google Sheets add-on and contribute:
\begin{enumerate}\itemsep0pt
\item \textbf{\sheetfault}, a seeded benchmark of workbooks, requests and \nBenchClasses{} classes of injected faults, labelled by an independent spreadsheet engine (not by the verifier under test) as loud or silent, with semantics-preserving variants as negatives.
\item \textbf{\gc}, a parser-based verifier with six workbook-grounded layers and structured repair feedback, replacing a regular-expression verifier, together with an account of what a workbook can and cannot decide about a formula (Section~\ref{sec:levels}).
\item \textbf{An evaluation whose labels and data are independent of the verifier}: real formulas from three disjoint Sheetpedia sets, the last two scored once with the verifier frozen; the production agent run with two local and three hosted models, with several sampling seeds for the local ones, and scored by execution; a model-critic baseline; and retrieval and cost experiments.
\item \textbf{Findings, including negative ones.} On \sheetfault{} the shipped verifier blocks \eonevonerejectR\% of faults and \gc{} \eonevtworejectR\% at \eonevtworejectFPR\% false positives. Because the benchmark was written with the verifier, we test on real formulas: the two verifiers reject about the same share there, rejections sit in a few workbooks, and the real data exposed a class of defects, deleted references, that the benchmark did not contain and \gc{} did not check. In end-to-end runs its feedback helps an 8B model more than resampling or a generic retry message (\sdADvtwofbresample{} points over resampling, in each of \sdASeeds{} seeds), but barely more than the shipped verifier's feedback (\sdADvtwofbvonefb{} points), and a 3B model gains only \sdBDvtwofbresample{}; hosted models make so few errors that the verifier is idle. Repair has a cost accuracy hides: it turns visible failures into invisible ones, raising the plausible-but-wrong formulas written from \hmAAsIsSilent{} to \hmARepVtwoSilent{} of \hmAN{} tasks for the 8B model, whereas blocking without repair does not.
\item \textbf{A secondary component}: a parse-and-classify URL validator for the fetch tool, which raises the block rate on an SSRF corpus from \ssrfvoneRecall\% to \ssrfvtwoRecall\% (Section~\ref{sec:e3}).
\end{enumerate}
